\section*{अध्याय \ref{ch:b1:diagenesis-basins} --- डायाजेनेसिस और अवसादी द्रोणियाँ}

\begin{solution}{exo:b1:diagenesis-basins:1}
बलुआ पत्थर $0.49\,e^{-2000/3704} = 0.49\times0.583 = \num{0.286}$; शेल
$0.63\,e^{-2000/1960} = 0.63\times0.360 = \num{0.227}$। छिद्रता आधी होती है जब
$e^{-z/\lambda} = 1/2$ हो, $z = \lambda\ln 2$ पर: बलुआ पत्थर के लिए \qty{2567}{m},
शेल के लिए \qty{1359}{m}।
\end{solution}

\begin{solution}{exo:b1:diagenesis-basins:2}
$\Delta A = \Delta E + S = 30 + 50 = \qty{80}{m}$। इसमें से \qty{60}{m} तलछट से
भरता है, इसलिए $\Delta W = 80 - 60 = \qty{20}{m}$: शेल्फ़ \qty{20}{m} गहरा होता है।
\end{solution}

\begin{solution}{exo:b1:diagenesis-basins:3}
$h = Y(\rho_m - \rho_w)/(\rho_m - \rho_s)$: \qty{2300}{kg.m^{-3}} के साथ,
$1.0\times2270/1000 = \qty{2.27}{km}$; \qty{2500}{kg.m^{-3}} के साथ,
$1.0\times2270/800 = \qty{2.84}{km}$। तलछट का हर मीटर स्तंभ में मैंटल के एक मीटर की
जगह लेता है; तलछट जितना सघन, मैंटल की तुलना में उसका द्रव्यमान-घाटा उतना छोटा,
इसलिए पानी से भरे छेद के समान घाटा भरने के लिए उतना ही अधिक चाहिए।
\end{solution}

\begin{solution}{exo:b1:diagenesis-basins:4}
(क)~कैल्साइट $100.09/2.71 = \qty{36.93}{cm^3.mol^{-1}}$; डोलोमाइट $184.40/2.84
= \qty{64.93}{cm^3.mol^{-1}}$। (ख)~दो मोल कैल्साइट घेरते हैं \qty{73.86}{cm^3}, एक
मोल डोलोमाइट \qty{64.93}{cm^3}: ठोस $1 - 64.93/73.86 = \qty{12.1}{\%}$ सिकुड़ता है।
(ग)~एक घन मीटर में \qty{0.90}{m^3} कैल्साइट $0.90\times(1 - 0.121) = \qty{0.791}{m^3}$
डोलोमाइट बन जाता है; छिद्रता $1 - 0.791 = \num{0.21}$ हो जाती है।
\end{solution}

\begin{solution}{exo:b1:diagenesis-basins:5}
(क)~$s = 1000 - 0.63\times1960\times(1 - e^{-1000/1960}) = \qty{506.5}{m}$; औसत
छिद्रता $1 - 506.5/1000 = \num{0.493}$ है। (ख)~ऊपरी सिरा \qty{2000}{m} पर होने पर,
$s(2000, b) = 506.5$ हल करने से $b = \qty{2628.7}{m}$ मिलता है: परत
\qty{628.7}{m} मोटी है, औसत छिद्रता $1 - 506.5/628.7 = \num{0.194}$ के साथ।
(ग)~कण वही हैं, इसलिए खोया पानी खोई मोटाई के बराबर है:
$1000 - 628.7 = \qty{371}{m^3}$ पानी प्रति वर्ग मीटर।
\end{solution}

\begin{solution}{exo:b1:diagenesis-basins:6}
(क)~$s = 300 - 0.49\times3704\times(e^{-2500/3704} - e^{-2800/3704}) =
\qty{228.10}{m}$। (ख)~$b_0 = 300$ से: $s(0, 300) = 158.80$, $1 - \varphi(300) =
0.5481$, $b_1 = 300 + (228.10 - 158.80)/0.5481 = \qty{426.44}{m}$; फिर
$s(0, 426.44) = 229.07$, $1 - \varphi = 0.5633$, $b_2 = 426.44 - 0.96/0.5633 =
\qty{424.73}{m}$; तीसरा चरण एक मिलीमीटर से कम बदलता है। बलुआ पत्थर सतह पर
\qty{424.7}{m} मोटा था। (ग)~$424.7/300 = 1.42$: दफ़न ने उसे \num{1.42} गुना पतला
किया है।
\end{solution}

\begin{solution}{exo:b1:diagenesis-basins:7}
(क)~बर्फ़ में $52\times10^6\times916.7/1000 = \qty{4.77e7}{km^3}$ पानी है; इसे
\qty{361.84e6}{km^2} पर फैलाने पर यह $\qty{0.1317}{km} = \qty{132}{m}$ समुद्र-स्तर है।
(ख)~तटरेखाएँ \qty{134}{m} देती हैं, \qty{2}{m} अधिक। (ग)~$4.767\times
10^7/0.134 = \qty{355.7e6}{km^2}$, आज से \qty{1.7}{\%} कम। जैसे-जैसे समुद्र-स्तर
गिरा, महाद्वीपीय शेल्फ़ उभर आए, इसलिए हिमनद महासागर छोटा था, और उससे लिया गया
उतना ही पानी उसे और अधिक नीचे गिराता था।
\end{solution}

\begin{solution}{exo:b1:diagenesis-basins:8}
(क)~$(-9.1 - (-98.1))/8 = \qty{11.1}{m}$ प्रति हज़ार वर्ष, तटरेखाओं से पाए
\qty{12}{m} के निकट। (ख)~\qty{40}{mm} प्रति वर्ष, यानी \qty{40}{m} प्रति हज़ार
वर्ष, $40/12 = 3.3$ गुना औसत है; पाँच सौ वर्षों में, $0.040\times500 = \qty{20}{m}$।
\end{solution}

\begin{solution}{exo:b1:diagenesis-basins:9}
(क)~$S = 50\times2 = \qty{100}{m}$ और $\Delta E = \qty{-40}{m}$: $\Delta A = \qty{60}{m}$,
सारा तलछट से भरा, क्योंकि गहराई अपरिवर्तित है। (ख)~\num{20000} वर्षों में
आधार-तल केवल $50\times0.02 = \qty{1}{m}$ धँसता है: \omterm{def:b1:diagenesis-basins:eustasy}{सापेक्ष समुद्र-स्तर}
$120 - 1 = \qty{119}{m}$ गिरता है। समुद्र-तल, जो समुद्र-स्तर से \qty{20}{m} नीचे था,
अब $119 - 20 = \qty{99}{m}$ ऊपर समाप्त होता है: पूरा शेल्फ़ उभर आता है।
\end{solution}

\begin{solution}{exo:b1:diagenesis-basins:10}
(क)~एक पतली परत के लिए छिद्रता उसके आर-पार लगभग स्थिर है: उसकी ठोस मोटाई
गहराई पर $s \approx h(1 - \varphi(z))$ है और सतह पर $s \approx h_0(1 -
\varphi_0)$ है। दोनों को बराबर करने पर $h_0 \approx h(1 -
\varphi(z))/(1 - \varphi_0)$ मिलता है। (ख)~$\varphi(3000) = 0.63\,e^{-3000/1960} =
0.136$, इसलिए $h_0 \approx 10\times0.864/0.37 = \qty{23.3}{m}$, जबकि सटीक
\qty{23.1}{m}। (ग)~सतह पर लौटाई गई परत की मोटाई \qty{23}{m} है, और अब वह सतह के
पास छिद्रता के क्षय की तुलना में पतली नहीं: उसकी छिद्रता ऊपरी सिरे पर
$\varphi_0$ से निचले सिरे पर कम तक गिरती है, इसलिए उसके कण सारी परत में
$\varphi_0$ से परिकलित मोटाई से कुछ कम में समाते हैं।
\end{solution}

\begin{solution}{exo:b1:diagenesis-basins:11}
(क)~$s = 3000 - 0.63\times1960\times(1 - e^{-3000/1960}) = \qty{2032}{m}$; औसत
छिद्रता $1 - 2032/3000 = \num{0.323}$ है और औसत घनत्व
$0.323\times1030 + 0.677\times2700 = \qty{2161}{kg.m^{-3}}$। (ख)~गुणक
$(3300 - 1030)/(3300 - 2161) = \num{1.99}$ है; भराव एक जल-भारित \omterm{def:b1:diagenesis-basins:subsidence}{अवधसान} दर्शाता है,
$3000/1.99 = \qty{1505}{m}$। (ग)~कणों के घनत्व से गुणक $2270/600 = \num{3.78}$ होगा
और \omterm{def:b1:diagenesis-basins:subsidence}{अवधसान} केवल \qty{793}{m}। यह गलत होगा, क्योंकि भराव का एक-तिहाई भाग पानी है:
तलछट अपने कणों से कहीं हल्की है, और द्रव्यमान-संतुलन में केवल संतृप्त तलछट का
घनत्व गिना जाता है।
\end{solution}

\begin{solution}{exo:b1:diagenesis-basins:12}
(क)~घन में \qty{0.90}{m^3} $= \qty{9.0e5}{cm^3}$ कैल्साइट है, यानी
$9.0\times10^5/36.93 = \num{24370}$ मोल। (ख)~हर दो मोल कैल्साइट पर एक मैग्नीशियम:
\num{12180} मोल। (ग)~समुद्री जल का एक किलोग्राम $35.165\times0.0365 =
\qty{1.284}{g}$ मैग्नीशियम रखता है, $1.284/24.305 = \qty{0.0528}{mol}$; अभिक्रिया को
$12180/0.0528 = \qty{2.31e5}{kg}$ समुद्री जल चाहिए, यानी $\qty{224}{m^3}$ ---
चट्टान के आयतन से दो सौ गुना से अधिक, भले ही मैग्नीशियम का हर आयन ले लिया जाए।
(घ)~डोलोमाइट बड़े पैमाने पर केवल वहीं बन सकता है जहाँ बहुत बड़े आयतन का पानी
लंबे समय तक चट्टान से बहे, या जहाँ पानी समुद्री जल से अधिक मैग्नीशियम वाला
ब्राइन हो: उदाहरण के लिए डूबते ब्राइन से धुले मंचों में, न कि ऐसे तलछट में जो
वही छिद्र-जल रखता है।
\end{solution}

\begin{solution}{pb:b1:diagenesis-basins:1}
\textbf{1.} बलुआ पत्थर: $\varphi(0) = \num{0.490}$, $\varphi(800) =
0.49\,e^{-800/3704} = \num{0.395}$। शेल: $\varphi(800) = \num{0.419}$,
$\varphi(2600) = \num{0.167}$। चॉक: $\varphi(2600) = \num{0.110}$,
$\varphi(3200) = \num{0.072}$।

\textbf{2.} \cref{prop:b1:diagenesis-basins:solid} से $s$ और $\bar\varphi
= 1 - s/h$ (प्रश्न~7): बलुआ पत्थर \num{0.441}, शेल \num{0.274}, चॉक \num{0.090}।

\textbf{3.} $h\bar\varphi$: बलुआ पत्थर \qty{353}{m^3}, शेल \qty{493}{m^3}, चॉक
\qty{54}{m^3} पानी प्रति वर्ग मीटर।

\textbf{4.} $\bar\varphi\rho_w + (1 - \bar\varphi)\rho_g$: बलुआ पत्थर
$0.441\times1030 + 0.559\times2650 = \qty{1936}{kg.m^{-3}}$, शेल
\qty{2242}{kg.m^{-3}}, चॉक \qty{2559}{kg.m^{-3}}। स्तंभ का भार
$1936\times800 + 2242\times1800 + 2559\times600 = \qty{7.12e6}{kg}$ प्रति वर्ग मीटर
है: औसत घनत्व \qty{2225}{kg.m^{-3}}।

\textbf{5.} शैलभार $2225\times9.81\times3200 = \qty{69.9}{MPa}$; जल-स्थैतिक
$1030\times9.81\times3200 = \qty{32.3}{MPa}$; \omterm{def:b1:stress-strain-rheology:lithostatic}{प्रभावी प्रतिबल}
$69.9 - 32.3 = \qty{37.5}{MPa}$।

\textbf{6.} $0.49\,e^{-z/3704} = 0.63\,e^{-z/1960}$ से $z = \ln(0.63/0.49)
/(1/1960 - 1/3704) = \qty{1046}{m}$। इससे नीचे बलुआ पत्थर अधिक छिद्रमय है: शेल गीला
शुरू करता है, पर अपना पानी कम दूरी में खोता है।

\textbf{7.} बलुआ पत्थर $800 - 0.49\times3704\times(1 - e^{-800/3704}) =
\qty{447.4}{m}$; शेल \qty{1306.7}{m}; चॉक \qty{546.0}{m}।

\textbf{8.} \qty{2300.2}{m} कण; पानी स्तंभ का $1 - 2300.2/3200 =
\qty{28}{\%}$ बनाता है।

\textbf{9.} बलुआ पत्थर $447.4/15 = \qty{29.8}{m}$, शेल $1306.7/40 =
\qty{32.7}{m}$, चॉक $546.0/20 = \qty{27.3}{m}$ ठोस प्रति दस लाख वर्ष।

\textbf{10.} शेल की आपूर्ति सबसे तेज़ थी, चॉक की सबसे धीमी --- पर तीनों में अंतर
केवल लगभग एक-पाँचवें का है, जो आज की उनकी मोटाइयाँ सुझाती हैं उससे कहीं कम।

\textbf{11.} \qty{800}{m}: उसका ऊपरी सिरा आज सतह पर है, जहाँ उसका निक्षेपण अभी-अभी
समाप्त हुआ, इसलिए निक्षेपण के समय उसकी मोटाई आज की मोटाई है।

\textbf{12.} शेल के लिए $s(0, b) = 1306.7$ को \cref{met:b1:diagenesis-basins:decompaction}
से हल करने पर $b = \qty{2124}{m}$ मिलता है: शेल \qty{15}{\Ma} पर \qty{2124}{m} मोटा था।

\textbf{13.} चॉक के लिए $s(0, b) = 546.0$ \qty{55}{\Ma} पर \qty{1071}{m} देता है।

\textbf{14.} ऊपरी सिरा \qty{2124}{m} पर होने पर, $s(2124, b) = 546.0$ से $b =
\qty{2748}{m}$ मिलता है: चॉक \qty{624}{m} मोटा था।

\textbf{15.} \qty{55}{\Ma} पर \qty{1071}{m}; $2124 + 624 = \qty{2748}{m}$ \qty{15}{\Ma} पर।

\textbf{16.} आज की मोटाइयों से: बलुआ पत्थर $800/15 = \qty{53.3}{m}$, शेल
$1800/40 = \qty{45.0}{m}$, चॉक $600/20 = \qty{30.0}{m}$ प्रति दस लाख वर्ष। निक्षेपण
के समय की मोटाइयों से: \qty{53.3}{m}, $2124/40 =
\qty{53.1}{m}$ और $1071/20 = \qty{53.6}{m}$। दूसरा समुच्चय सार्थक है --- वह दर जिस पर
समुद्र-तल बना --- और वह एक स्थिर दर दिखाता है जिसे आज की मोटाइयाँ छिपाती हैं।

\textbf{17.} $1 - 1800/2124 = \qty{15}{\%}$।

\textbf{18.} पानी की गहराई हर अंतराल के आरंभ और अंत पर एक ही है, इसलिए बना
\omterm{def:b1:diagenesis-basins:accommodation}{आवास-स्थान} आधार-तल के ऊपर स्तंभ की वृद्धि है: \qty{75}{\Ma} से \qty{55}{\Ma} तक
\qty{1071}{m}, \qty{55}{\Ma} से \qty{15}{\Ma} तक $2748 - 1071 = \qty{1677}{m}$, और
\qty{15}{\Ma} से $3200 - 2748 = \qty{452}{m}$।

\textbf{19.} $S = \Delta A - \Delta E$: \qty{55}{\Ma} से \qty{15}{\Ma} तक
$1677 - (-60) = \qty{1737}{m}$; उसके बाद से $452 - (-30) = \qty{482}{m}$।

\textbf{20.} $1737/40 = \qty{43}{m}$ और $482/15 = \qty{32}{m}$ प्रति दस लाख वर्ष।

\textbf{21.} $60/1737 = \qty{3.5}{\%}$ और $30/482 = \qty{6.2}{\%}$। दोनों अंतरालों में
समुद्र-स्तर गिरा, इसलिए \omterm{def:b1:diagenesis-basins:eustasy}{यूस्टेसी} ने भराव में बाधा डाली: आधार-तल को उतना ही अधिक
धँसना पड़ा ताकि वही जगह बने। लगभग सारा \omterm{def:b1:diagenesis-basins:accommodation}{आवास-स्थान} \omterm{def:b1:diagenesis-basins:subsidence}{अवधसान} ने बनाया।

\textbf{22.} शेल और चॉक \qty{15}{\Ma} पर \qty{2748}{m} मोटे थे और आज \qty{2400}{m} हैं:
उनके \omterm{def:b1:diagenesis-basins:compaction}{संपीड़न} ने \qty{348}{m} जगह बनाई, बलुआ पत्थर द्वारा भरे गए \qty{452}{m} का
\qty{77}{\%}।

\textbf{23.} हर इकाई को उसके बाद के दफ़न ने पतला किया है, उस मात्रा से जो उसके
शैल-प्रकार और इस पर निर्भर करती है कि वह कितना गहरा गया; हर अंतराल में बनी जगह
उन मोटाइयों से मापी जानी चाहिए जो इकाइयों की तब थीं, जो केवल \omterm{def:b1:diagenesis-basins:compaction}{डिकंपैक्शन} देता
है। और तब भी इस प्रकार मिला \omterm{def:b1:diagenesis-basins:subsidence}{अवधसान} तलछट के भार से हुए धँसाव को शामिल करता है,
जिसे \omterm{def:b1:diagenesis-basins:subsidence}{विवर्तनिक अवधसान} तक पहुँचने के लिए हटाना होगा।

\textbf{24.} \qty{2124}{m}, लगभग \qty{2.1}{km} --- कुएँ में आज मिली \qty{1800}{m} से
\qty{324}{m} अधिक।
\end{solution}
